\documentclass{article}

\usepackage[boxed,commentsnumbered]{algorithm2e}

\newcommand{\algocfversion}{release 5.2}
\title{BookML test: algorithm2e.sty --- package for algorithms\\ {\large\algocfversion}}
\author{(c) 1995-1997 Christophe Fiorio, Tu-Berlin, Germany\\
(c) 1998-2017 Christophe Fiorio, LIRMM, Montpellier University, France} \date{July 18 2017}

\usepackage{bookml/bookml}
\usepackage{framed}
\AtBeginDocument{\begin{framed}
  \begin{description}
    \item[BookML] \BookMLversion
    \item[\LaTeXML] \LaTeXMLversion{} (\csname bml@generator@engine\endcsname)
    \item[\LaTeX] \fmtversion
  \end{description}
\end{framed}}


\begin{document}

\begin{algorithm}[H]
  \SetAlgoLined
  \KwData{this text}
  \KwResult{how to write algorithm with \LaTeX2e }

  initialization\;
  \While{not at end of this document}{
    read current section\;
    \eIf{understand}{
      go to next section\;
      current section becomes this one\;
      }{
      go back to the beginning of current section\;
      }
    }
  \caption{How to write algorithms}
\end{algorithm}

\RestyleAlgo{boxed}\LinesNumbered
\IncMargin{1em}
\begin{algorithm}
  \SetKwData{Left}{left}\SetKwData{This}{this}\SetKwData{Up}{up}
  \SetKwFunction{Union}{Union}\SetKwFunction{FindCompress}{FindCompress}
  \SetKwInOut{Input}{input}\SetKwInOut{Output}{output}

  \Input{A bitmap $Im$ of size $w\times l$}
  \Output{A partition of the bitmap}
  \BlankLine
  \emph{special treatment of the first line}\;
  \For{$i\leftarrow 2$ \KwTo $l$}{
    \emph{special treatment of the first element of line $i$}\;
    \For{$j\leftarrow 2$ \KwTo $w$}{\label{forins}
      \Left$\leftarrow$ \FindCompress{$Im[i,j-1]$}\;
      \Up$\leftarrow$ \FindCompress{$Im[i-1,]$}\;
      \This$\leftarrow$ \FindCompress{$Im[i,j]$}\;
      \If(\tcp*[h]{O(\Left,\This)==1}){\Left compatible with \This}{\label{lt}
        \lIf{\Left $<$ \This}{\Union{\Left,\This}}
        \lElse{\Union{\This,\Left}}
      }
      \If(\tcp*[f]{O(\Up,\This)==1}){\Up compatible with \This}{\label{ut}
        \lIf{\Up $<$ \This}{\Union{\Up,\This}}
        \tcp{\This is put under \Up to keep tree as flat as possible}\label{cmt}
        \lElse{\Union{\This,\Up}}\tcp*[h]{\This linked to \Up}\label{lelse}
      }
    }
    \lForEach{element $e$ of the line $i$}{\FindCompress{p}}
  }
  \caption{disjoint decomposition}\label{algo_disjdecomp}
\end{algorithm}\DecMargin{1em}

{\RestyleAlgo{algoruled}\SetAlgoVlined\LinesNotNumbered
\begin{algorithm}
\DontPrintSemicolon
\KwData{$G=(X,U)$ such that $G^{tc}$ is an order.}
\KwResult{$G'=(X,V)$ with $V\subseteq U$ such that $G'^{tc}$ is an
interval order.}
\Begin{
  $V \longleftarrow U$\;
  $S \longleftarrow \emptyset$\;
  \For{$x\in X$}{
    $NbSuccInS(x) \longleftarrow 0$\;
    $NbPredInMin(x) \longleftarrow 0$\;
    $NbPredNotInMin(x) \longleftarrow |ImPred(x)|$\;
    }
  \For{$x \in X$}{
    \If{$NbPredInMin(x) = 0$ {\bf and} $NbPredNotInMin(x) = 0$}{
      $AppendToMin(x)$}
    }
    \nl\While{$S \neq \emptyset$}{\label{InRes1}
    \nlset{REM} remove $x$ from the list of $T$ of maximal index\;\label{InResR}
    \lnl{InRes2}\While{$|S \cap  ImSucc(x)| \neq |S|$}{
      \For{$ y \in  S-ImSucc(x)$}{
        \{ remove from $V$ all the arcs $zy$ : \}\;
        \For{$z \in  ImPred(y) \cap  Min$}{
          remove the arc $zy$ from $V$\;
          $NbSuccInS(z) \longleftarrow NbSuccInS(z) - 1$\;
          move $z$ in $T$ to the list preceding its present list\;
          \{i.e. If $z \in T[k]$, move $z$ from $T[k]$ to
           $T[k-1]$\}\;
          }
        $NbPredInMin(y) \longleftarrow 0$\;
        $NbPredNotInMin(y) \longleftarrow 0$\;
        $S \longleftarrow S - \{y\}$\;
        $AppendToMin(y)$\;
        }
      }
    $RemoveFromMin(x)$\;
    }
  }
\caption{IntervalRestriction\label{IR}}
\end{algorithm}}

\def\theAlgoLine{\arabic{AlgoLine}}
\SetKwProg{Fn}{Function}{}{end}\SetKwFunction{FRecurs}{FnRecursive}%
\newcommand{\forcond}{$i=0$ \KwTo $n$}
\begin{algorithm}\SetAlgoLongEnd
  \caption{Generic example with most classical expressions derived in pseudo-code}%
  \label{algo:pseudocode}
\Fn(\tcc*[h]{algorithm as a recursive function}){\FRecurs{some args}}{
  \KwData{Some input data\\these inputs can be displayed on several lines and one
    input can be wider than line's width.}
  \KwResult{Same for output data}
  \tcc{this is a comment to tell you that we will now really start code}
  \If(\tcc*[h]{a simple if but with a comment on the same line}){this is true}{
    we do that, else nothing\;
    \tcc{we will include other if so you can see this is possible}
    \eIf{we agree that}{
      we do that\;
    }{
      else we will do a more complicated if using else if\;
      \uIf{this first condition is true}{
        we do that\;
      }
      \uElseIf{this other condition is true}{
        this is done\tcc*[r]{else if}
      }
      \Else{
        in other case, we do this\tcc*[r]{else}
      }
    }
  }
  \tcc{now loops}
  \For{\forcond}{
    a for loop\;
  }
  \While{$i<n$}{
    a while loop including a repeat--until loop\;
    \Repeat{this end condition}{
      do this things\;
    }
  }
  They are many other possibilities and customization possible that you have to
  discover by reading the documentation.
}
\end{algorithm}

\end{document}
